% Physical pages: 182--187 (printed pages 159--164).
% Figures: none.
% Uncertainties: none.

% Continuation of Section 3.7 from physical page 181.
The smallness of the real part of the forward scattering amplitude at high
energy is confirmed by experiment.

So far we have not said anything about whether the interaction radius $R_n$
itself may depend on energy. As a very crude guess, we may take $R_n$ as the
distance at which the factor $\exp(-\mu r)$ in the Yukawa potential \eqref{eq:1.2.74}
takes a value proportional to some unknown power of $E$, in which case $R_n$
goes as $\log E$ for $E\to\infty$, and the cross-sections go as $(\log E)^2$.
As it happens, it has been rigorously shown~\hyperref[ref:3.26]{[26]} on the basis of very general
assumptions that the total cross-section can grow no faster than $(\log E)^2$
for $E\to\infty$, and in fact the observed proton--proton total cross-section
rises something like $(\log E)^2$ at high energies, so this rough picture of
high-energy scattering does seem to have some correspondence with reality.

\subsection[Resonances]{Resonances\texorpdfstring{\protect\footnotemark}{}}
\footnotetext{This section lies somewhat out of the book's main line of
development, and may be omitted in a first reading.}

It often happens that the particles participating in a multi-particle
collision can form an intermediate state consisting of a \emph{single}
unstable particle $R$, that eventually decays into the particles observed as
the final state. If the total decay rate of $R$ is small the cross-section
exhibits a rapid variation (usually a peak), known as a \emph{resonance}, at
the energy of the intermediate state $R$.

We shall see that the behavior of the cross-sections near a resonance is
pretty much prescribed by the unitarity condition alone, which is a good
thing since there are a number of very different mechanisms that can produce
a nearly stable state:

\begin{list}{(\alph{enumi})}{\usecounter{enumi}\setlength{\leftmargin}{2.3em}}
\item The simplest possibility is that the Hamiltonian can be decomposed into
two terms, a `strong' Hamiltonian $H_0+H_s$, which has the particle $R$ as an
eigenstate, plus a weak perturbation $H_w$, which allows $R$ to decay into
various states, including the initial and final states $\alpha,\beta$ of our
collision process. For instance there is a neutral particle, the $Z^0$, with
$j=1$ and mass 91 GeV, that would be stable in the absence of the electroweak
interactions. These interactions allow the $Z^0$ to decay into
electron--positron pairs, muon--antimuon pairs, etc., but with a total decay
rate that is much less than the $Z^0$ mass. In 1989 the $Z^0$ particle was
seen as a resonance\footnote{Incidentally, this example shows that a resonant
state only needs to decay \emph{relatively} slowly; the $Z^0$ lifetime is
$2.6\cdot10^{-25}$ seconds, which is not long enough for a $Z^0$ travelling
near the speed of light to cross an atomic nucleus. What is important is that
the decay rate is 36 times smaller than the rate $M_Z/\hbar$ of oscillation of
the $Z^0$ wave function in its rest frame.} in electron--positron collisions
at CERN and Stanford, in the reactions
$e^++e^-\to Z^0\to e^++e^-$,
$e^++e^-\to Z^0\to\mu^++\mu^-$, etc.

\item In some cases a particle is long-lived because there is a potential
barrier that nearly prevents its constituents from escaping. The classic
example is nuclear alpha decay: it may be energetically possible for a
nucleus to emit an alpha particle (a He$^4$ nucleus) but the strong
electrostatic repulsion between the alpha particle and the nucleus creates a
barrier region around the daughter nucleus, which the alpha particle is
classically forbidden to enter. The decay then can proceed only by quantum
mechanical barrier penetration, and is exponentially slow. Such an unstable
state shows up as a resonance in the scattering of the alpha particle on the
daughter nucleus. For instance, the lowest-energy state of the Be$^8$ nucleus
is unstable against decay into two alpha particles, and is seen as a
resonance in He$^4$--He$^4$ scattering. (In addition to Coulomb barriers,
there are also centrifugal barriers that help lengthen the life of alpha-,
beta-, and gamma-unstable nuclei of high spin.)

\item It is possible for complicated systems to be nearly unstable for
statistical reasons, without the presence of any potential barriers or weak
interactions. For instance, an excited state of a heavy nucleus may be able to
decay only if, through a statistical fluctuation, a large part of its energy
is concentrated on a single neutron. This state will then show up as a
resonance in the scattering of a neutron on the daughter nucleus.
\end{list}

These mechanisms for producing long-lived states are so different that it is
truly fortunate that most of the properties of resonances follow from
unitarity alone, without regard to the dynamical mechanism that produces the
resonance.

First, let's consider the energy dependence of the matrix element for a
reaction near a resonance. A wave packet
$\int d\alpha\,g(\alpha)\ket{\alpha}_{\mathrm{in}}\exp(-iE_\alpha t)$ of `in'
states has a time-dependence given by \eqref{eq:3.1.19}
\[
  \int d\alpha\,g(\alpha)\ket{\alpha}_{\mathrm{in}}e^{-iE_\alpha t}
  =\int d\alpha\,g(\alpha)\ket{\alpha}_0e^{-iE_\alpha t}
  +\int d\beta\,\ket{\beta}_0\int d\alpha\,
   \frac{e^{-iE_\alpha t}g(\alpha)T^+_{\beta\alpha}}
        {E_\alpha-E_\beta+i\epsilon}.
\]
As mentioned in Section 3.1, a pole in the function $T^+_{\beta\alpha}$ in
the lower-half complex $E_\alpha$ plane would make a contribution to the
second term that decays exponentially as $t\to\infty$. Specifically, a pole
at $E_\alpha=E_R-i\Gamma/2$ yields a term in the \emph{amplitude} that behaves
like $\exp(-iE_Rt-\Gamma t/2)$, so it corresponds to a state whose
\emph{probability} decays like $\exp(-\Gamma t)$. We conclude then that a
long-lived state of energy $E_R$ with a slow decay rate $\Gamma$ produces a
term in the scattering amplitude that varies as
\begin{equation}
  T^+_{\beta\alpha}\sim(E_\alpha-E_R+i\Gamma/2)^{-1}+\text{constant}.
  \tag{3.8.1}\label{eq:3.8.1}
\end{equation}

To go further, it will be convenient to adopt as a basis the orthonormal
discrete multi-particle states $\ket{\mathbf p,E,N}$ discussed in the previous
section; $\mathbf p$ and $E$ are the total momentum and energy, and $N$ is an
index that takes only discrete (though infinitely many) values. In this basis,
the $S$-matrix may be written
\begin{equation}
  S_{\mathbf p'E'N',\mathbf pEN}
  =\delta^3(\mathbf p'-\mathbf p)\delta(E'-E)S_{N'N}(\mathbf p,E).
  \tag{3.8.2}\label{eq:3.8.2}
\end{equation}
Near a resonance, we expect the center-of-mass frame amplitude
$S(\mathbf0,E)\equiv\mathcal S(E)$ to have the form
\begin{equation}
  \mathcal S_{N'N}(E)\equiv S_{N'N}(\mathbf0,E)
  =\mathcal S_{0N'N}+\frac{\mathcal R_{N'N}}{E-E_R+i\Gamma/2},
  \tag{3.8.3}\label{eq:3.8.3}
\end{equation}
where $\mathcal S_0$ and $\mathcal R$ are approximately constant at least over
the relatively small range of energies $|E-E_R|\lesssim\Gamma$.

In this basis, the unitarity of the $S$-matrix is an ordinary matrix equation
\begin{equation}
  \mathcal S(E)^\dagger\mathcal S(E)=\mathbf{1}.
  \tag{3.8.4}\label{eq:3.8.4}
\end{equation}
Applied to Eq. \eqref{eq:3.8.3}, this tells us that the non-resonant background
$S$-matrix is unitary
\begin{equation}
  \mathcal S_0^\dagger\mathcal S_0=\mathbf{1},
  \tag{3.8.5}\label{eq:3.8.5}
\end{equation}
and also that the residue matrix $\mathcal R$ satisfies the two conditions
\begin{equation}
  \mathcal S_0^\dagger\mathcal R-\mathcal R^\dagger\mathcal S_0=0,
  \tag{3.8.6}\label{eq:3.8.6}
\end{equation}
\begin{equation}
  -\frac{i}{2}\Gamma\mathcal S_0^\dagger\mathcal R
  +\frac{i}{2}\Gamma\mathcal R^\dagger\mathcal S_0
  +\mathcal R^\dagger\mathcal R=0.
  \tag{3.8.7}\label{eq:3.8.7}
\end{equation}
These conditions can be put in a more transparent form by setting
\begin{equation}
  \mathcal R\equiv-i\Gamma\mathcal A\mathcal S_0.
  \tag{3.8.8}\label{eq:3.8.8}
\end{equation}
The unitarity conditions on the matrix $\mathcal A$ are then simply
\begin{equation}
  \mathcal A^\dagger=\mathcal A,
  \qquad \mathcal A^2=\mathcal A.
  \tag{3.8.9}\label{eq:3.8.9}
\end{equation}
Any such Hermitian idempotent matrix is called a \emph{projection matrix}.
Such matrices can always be expressed as a sum of dyads of orthonormal vectors
$u^{(r)}$:
\begin{equation}
  \mathcal A_{N'N}=\sum_r u_{N'}^{(r)}u_N^{(r)*},
  \qquad \sum_Nu_N^{(r)*}u_N^{(s)}=\delta_{rs}.
  \tag{3.8.10}\label{eq:3.8.10}
\end{equation}
The discrete part of the $S$-matrix is then
\begin{equation}
  \mathcal S_{N'N}(E)=\sum_{N''}
  \left[\delta_{N'N''}-i\frac{\Gamma}{E-E_R+i\Gamma/2}
  \sum_r u_{N'}^{(r)}u_{N''}^{(r)*}\right]\mathcal S_{0N''N}.
  \tag{3.8.11}\label{eq:3.8.11}
\end{equation}

Each term in the sum over $r$ can be thought of as arising from a different
resonant state, all these states having the same values for $E_R$ and
$\Gamma$.

What has this to do with rates and cross-sections? For simplicity let's now
ignore the non-resonant background scattering, setting
$\mathcal S_{0N'N}$ equal to $\delta_{N'N}$; we will come back to the more
general case a little later. Then for the two-body discrete center-of-mass
states described in the previous section, Eq. \eqref{eq:3.8.11} reads:
\begin{align}
  \mathcal S_{j'\sigma'\ell's'n',j\sigma\ell sn}(E)
  &=\delta_{j'j}\delta_{\sigma'\sigma}\delta_{\ell'\ell}
    \delta_{s's}\delta_{n'n} \notag\\
  &\quad-i\frac{\Gamma}{E-E_R+i\Gamma/2}
    \sum_r u_{j'\sigma'\ell's'n'}^{(r)}
           u_{j\sigma\ell sn}^{(r)*}.
  \tag{3.8.12}\label{eq:3.8.12}
\end{align}
In all cases the label $r$ will include an index $\sigma_R$ giving the
$z$-component of the total angular momentum of the resonant state; for a
resonant state of total angular momentum $j_R$, $\sigma_R$ takes $2j_R+1$
values. If there is no other degeneracy, then $r$ just labels the value of
$\sigma_R$, and
\begin{equation}
  u_{j\sigma\ell sn}^{(\sigma_R)}
  =\delta_{jj_R}\delta_{\sigma\sigma_R}u_{\ell sn},
  \tag{3.8.13}\label{eq:3.8.13}
\end{equation}
where $u_{\ell sn}$ are a set of complex amplitudes that (because of the
Wigner--Eckart theorem) are independent of $\sigma$. Now Eq. \eqref{eq:3.8.12} gives
the amplitude $S^j$ defined by Eq. \eqref{eq:3.7.7} as
\begin{equation}
  S^j_{\ell's'n',\ell sn}(E)
  =\delta_{\ell'\ell}\delta_{s's}\delta_{n'n}
  -i\delta_{jj_R}\frac{\Gamma}{E-E_R+i\Gamma/2}
   u_{\ell's'n'}u_{\ell sn}^*.
  \tag{3.8.14}\label{eq:3.8.14}
\end{equation}
Also, Eq. \eqref{eq:3.8.10} now reads
\begin{equation}
  \sum_{\ell sn}|u_{\ell sn}|^2+\cdots=1
  \tag{3.8.15}\label{eq:3.8.15}
\end{equation}
with the dots representing the positive contribution of any states containing
three or more particles. As we shall see, the quantities $|u_{\ell sn}|^2$
have the interpretation of branching ratios for the decay of the resonant
state into the various accessible two-body states.

Eq. \eqref{eq:3.7.12} now gives the total cross-section for all reactions in channel
$n$:
\begin{equation}
  \sigma_{\rm total}(n;E)
  =\frac{\pi(2j_R+1)}{k^2(2s_1+1)(2s_2+1)}
   \frac{\Gamma\Gamma_n}{(E-E_R)^2+\Gamma^2/4},
  \tag{3.8.16}\label{eq:3.8.16}
\end{equation}
where
\begin{equation}
  \Gamma_n\equiv\Gamma\sum_{\ell s}|u_{\ell sn}|^2.
  \tag{3.8.17}\label{eq:3.8.17}
\end{equation}
This is a version of the celebrated Breit--Wigner single-level formula~\hyperref[ref:3.27]{[27]}.
We can also use these results to calculate the cross-section for resonant
scattering from an initial two-body channel $n$ to a final two-body channel
$n'$. Using Eq. \eqref{eq:3.8.14} in Eq. \eqref{eq:3.7.10} gives
\begin{equation}
  \sigma(n\to n';E)
  =\frac{\pi(2j_R+1)}{k^2(2s_1+1)(2s_2+1)}
   \frac{\Gamma_n\Gamma_{n'}}{(E-E_R)^2+\Gamma^2/4}.
  \tag{3.8.18}\label{eq:3.8.18}
\end{equation}
This shows that the probabilities that the resonant state will decay into any
one of the final two-body channels $n'$ are proportional to the
$\Gamma_{n'}$. According to Eq. \eqref{eq:3.8.15}, the sum of the $\Gamma_n$
(including contributions from final states containing three or more
particles) is just equal to the total decay rate $\Gamma$, so we can conclude
that $\Gamma_n$ is just the rate for the decay of the resonant state into
channel $n$.

We see in Eqs. \eqref{eq:3.8.16} and \eqref{eq:3.8.18} the characteristic resonant peak at
energy $E_R$, with a width (the full width at half maximum) equal to the decay
rate $\Gamma$. (The individual $\Gamma_n$ are often called partial widths.)
Since $\Gamma_n\leq\Gamma$, the total cross-section at the peak of the
resonance is roughly bounded by one square wavelength, $(2\pi/k)^2$. This
rule, that cross-sections at a single resonance are roughly bounded by a
square wavelength, is universally applicable even in classical physics
(where energy conservation plays the role played here by unitarity), as for
instance in the resonant interaction of sound waves with bubbles in the sea,
or gravitational waves with gravitational wave antennae. (In the latter
case, the branching ratio for oscillations in any laboratory mass to lose
their energy through gravitational radiation is tiny, so the cross-section
even at a resonance peak is vastly less then a square wavelength~\hyperref[ref:3.28]{[28]}.)

Incidentally, it often happens that a resonance is detected, but energy
measurements are insufficiently precise to resolve its width. In this case,
what is measured experimentally is the integral of the cross-section over the
resonant peak. For the total cross-section \eqref{eq:3.8.14}, this is
\begin{equation}
  \int\sigma_{\rm total}(n;E)\,dE
  =\frac{2\pi^2(2j_R+1)\Gamma_n}
   {k_R^2(2s_1+1)(2s_2+1)}.
  \tag{3.8.19}\label{eq:3.8.19}
\end{equation}
Such experiments can reveal only the partial width for decay of the resonant
state into the initial particles, not the total width or branching ratios.

This formalism can also be applied when the resonant states with a given spin
$z$-component form a multiplet related by some symmetry group. For instance,
to the extent that isospin symmetry is respected, for a resonance of total
isospin $T_R$ the index $r$ labelling resonant states includes a specification
not only of the angular-momentum $z$-component $\sigma_R$, but also of an
isospin three-component $t_R$, taking values $-T_R,-T_R+1,\ldots,T_R$. In
this case there is no change in the above results for the total and partial
cross-sections, because each two-body channel $n$ has definite values $t_1,t_2$
of the isospin $z$-components of the two particles, and hence can only couple
to the resonant state with the single value $t_1+t_2$ for $t_R$. The partial
widths $\Gamma_n$ here depend on $t_1$ and $t_2$ only through factors
$C_{T_1T_2}(T_R,t_R;t_1,t_2)^2$.

The presence of a resonance shows itself in a characteristic behavior of
phase shifts near the resonance. Returning to the general formula \eqref{eq:3.8.11}
(but still taking $\mathcal S_0=\mathbf{1}$), we see from Eq. \eqref{eq:3.8.10} that
for each individual resonant state $r$, there is an eigenvector $u_N^{(r)}$
of $\mathcal S_{N'N}(E)$ with eigenvalue
\[
  \exp[2i\delta^{(r)}(E)]
  =1-i\frac{\Gamma}{E-E_R+i\Gamma/2}
\]
or, in other words,
\begin{equation}
  \tan\delta^{(r)}(E)=-\frac{\Gamma/2}{E-E_R}.
  \tag{3.8.20}\label{eq:3.8.20}
\end{equation}
We see that over an energy range of order $\Gamma$ centered around the
resonant energy, the `eigenphase' $\delta^{(r)}(E)$ jumps from a value
$\nu\pi$ (with $\nu$ a positive or negative integer) below the resonance to
$(\nu+1)\pi$ above it. However, in order to use this result to say something
about reaction rates, we need to know the eigenvectors $u_N^{(r)}$, which, in
general, have components with arbitrary numbers of particles with various
momenta, spin, and species.

These results are much more useful in those special cases when the particles
in a particular channel $N$ are forbidden (usually by conservation laws) from
making a transition to any other channel. With this assumption, it is not
difficult to include the effects of a non-resonant background scattering
matrix $\mathcal S_0$ in the general result \eqref{eq:3.8.11}. In order for
$\mathcal S_{N'N}$ to vanish for some particular $N$ and all $N'\neq N$, it
is necessary that the same is true of $\mathcal S_{0N'N}$, and also true of
$u_{N'}^{(r)}$ for any $r$ for which $u_N^{(r)}\neq0$. The unitarity
requirement \eqref{eq:3.8.5} then requires that for this $N$
\[
  \mathcal S_{0N'N}=\exp(2i\delta_{0N})\delta_{N'N}
\]
and Eq. \eqref{eq:3.8.10} requires that
\[
  u_N^{(r)*}u_N^{(s)}=\delta_{rs},
\]
so that there can be only one term $r$ in Eq. \eqref{eq:3.8.11} for which
$u_N^{(r)}\neq0$. In this case, Eq. \eqref{eq:3.8.11} gives
\[
  \mathcal S_{N'N}(E)=\delta_{N'N}
   \left[1-\frac{i\Gamma}{E-E_R+i\Gamma/2}\right]
   \exp(2i\delta_{0N})
  \equiv\delta_{N'N}\exp[2i\delta_N(E)]
\]
with total phase shift
\begin{equation}
  \delta_N(E)=\delta_{0N}-\arctan\!\left(\frac{\Gamma/2}{E-E_R}\right).
  \tag{3.8.21}\label{eq:3.8.21}
\end{equation}
We see that over an energy range of order $\Gamma$ centered around the
resonant energy $E_R$, the phase shift $\delta_N(E)$ jumps from a value
$\delta_{0N}$ below the
